\section{Introduction}

China's local government financing vehicle (LGFV) exit reform poses a
measurement problem before it poses an explanation problem. Official reports
record reductions in platform counts, financing balances, and local government
debt exposure, but those figures do not reveal what happened to the government
financing and public-project functions that made the platforms important.
Official exit is not institutional exit. The same compliance statement can
describe a firm that stopped quasi-fiscal financing, a firm that continued the
same function under new language, or a reorganization that moved the function
to another local state-owned entity.

This paper measures the institutional fate of official exit. The coding rule
begins with a documented formal event and then traces the post-event location
of financing, construction, land development, liabilities, fiscal support, and
state-asset control. It assigns one of four mutually exclusive outcomes.
Substantive exit means that the issuer no longer performs the relevant
government financing function. Nominal exit means that formal compliance
language appears while the same issuer continues that function. Functional
transfer means that the original issuer exits or is reorganized while the
function moves to another local state-owned entity. Liquidation records legal
termination and remains separate because termination alone does not establish
where the fiscal function went.

The framework adapts established concerns about organizational decoupling and
LGFV transformation to a source-auditable case measure. Prior work already
shows that formal policy can diverge from organizational practice, that list
exit can coexist with continued borrowing, and that platform transformation can
change organizational form without eliminating fiscal functions
\citep{bromley_powell_2012_decoupling,zhelyazkova_kaya_schrama_2016_compliance,
fang_lu_2024_classified_regulation,li_feng_hao_2022_lgfp_transformation,
feng_wu_zhang_2022_jiaxing}. The contribution here is not the general
formal-substantive distinction. It is the joint evidence design: verify the
formal event, trace the function after that event, assign a four-category fate,
and retain the source packet and alternative interpretation for review.

The current working-reference file contains ninety-four city-platform cases
assigned through Codex source-packet review on behalf of the project author.
They await independent human confirmation. The file records two substantive
exits, eighty-two nominal exits, ten
functional transfers, and no liquidation cases. These frequencies are not
national prevalence estimates because the cases were assembled to develop and
stress-test the measure. They nevertheless establish the empirical puzzle.
Identical no-government-financing language appears in cases where former
projects moved into fiscal channels, cases where the same issuer retained
public-project and fiscal-support functions, and cases where those functions
moved through a reorganized municipal platform system. All evidence identifiers
cited by the reference labels resolve against the document and source
inventories; ninety-three cases have complete local copies of the cited
documents, and every case has a recovery URL for each cited document.

I use the new outcome in an exploratory application to historically rooted
local capacity. LGFV exit can require three contemporary processes: fiscal
absorption of eligible obligations, bureaucratic coordination of assignments
and transfers, and project or state-asset governance after reorganization.
These processes are observable in budgets, debt registers, cross-agency
decisions, transfer records, contracts, and project-level control. Historical
elite density may be associated with the institutional environments that
support them, but it is not a direct measure of any one process and also
captures education and social capital \citep{chen_kung_ma_2020_keju}. The paper
therefore treats historical capacity as an explanatory application rather than
as part of the measurement contribution.

Among the eighty-four reference cases matched to Ming-Qing elite density,
simple specifications show a positive association between elite density and
movement away from nominal exit. The outcome contains only twelve such events,
however, and the complete-control subset contains seventy-eight observations.
It excludes six low-capacity nominal exits, and the initial adjusted design is
rank deficient and sensitive to individual case deletion. These diagnostics
do not overturn the descriptive pattern, but they rule out a strong claim that
historical capacity has a stable adjusted or causal effect. The empirical
application is retained to show how the measure can be used and what a stronger
design must solve.

Large language models assist source screening but do not define the outcome.
The expanded file contains 361 candidate disclosures and 203 one-sided
nominal-exit surrogate labels. Sixty-one surrogate issuers overlap the
working-reference file and all receive the same nominal-exit label. That
result is selected-overlap concordance, not a probability-sample estimate of
precision. No screen-negative issuer has a reference-label overlap, so recall,
calibration, and four-category error are not identified. Consistent with
design-based supervised learning, downstream correction requires a validation
sample with known inclusion probabilities and independently assigned human
labels \citep{egami2023dsl}. The present working labels support workflow
development but do not satisfy that final validation requirement. The paper
reports this boundary instead of
using the surrogate labels as if they were observed outcomes.

The paper makes three ordered contributions. First, it provides an auditable
measure of the institutional fate of LGFV exit. Second, it uses that measure to
identify descriptive variation that official counts conceal and to formulate a
restrained application to historical capacity. Third, it specifies a scalable
LLM-assisted workflow in which machine labels remain provisional, reference
labels remain tied to original documents, and unresolved measurement error is
carried into the inferential design. The following sections locate the measure
in the relevant literatures, describe the fiscal setting, develop the
institutional categories and contemporary processes, present the research and
coding design, and report the current evidence and its limits.
